% ============================================================
% ARIA-ICNABS Project — Paper v9
% A Diagnostic Framework for Detecting False Fractionalization
% in Physics-Informed Neural Networks
% Compile: pdflatex (2 runs for references)
% ============================================================
\documentclass[11pt]{article}

\usepackage[T1]{fontenc}
\usepackage[utf8]{inputenc}
\usepackage{amsmath, amssymb, amsthm}
\usepackage{graphicx}
\usepackage{geometry}
\usepackage{booktabs}
\usepackage{multirow}
\usepackage{enumitem}
\usepackage{hyperref}
\usepackage{xcolor}

\geometry{top=30mm, bottom=25mm, left=25mm, right=25mm}
\graphicspath{{./figures/}}

\newtheorem{proposition}{Proposition}
\newtheorem{definition}{Definition}
\newtheorem{remark}{Remark}

\title{\textbf{A Diagnostic Framework for Detecting False Fractionalization in Physics-Informed Neural Networks: \\ A Three-Way Comparison of Fractional Advection-Dispersion Residuals}}

\author{Maryam Jalili\thanks{Department of Mathematics, Neyshabur Branch, Islamic Azad University, Neyshabur, Iran. Email: mrmjlili@iau.ac.ir}}

\date{}

\begin{document}
\maketitle

% ============================================================
% ABSTRACT
% ============================================================
\begin{abstract}
\noindent
Physics-informed neural networks (PINNs) are increasingly deployed for fractional partial differential equations, yet the field lacks a systematic method to verify that a chosen residual discretization genuinely preserves the non-local memory of the underlying fractional operator. We report a case study on the Atangana--Baleanu--Caputo (ABC) derivative in a one-dimensional advection--dispersion equation, in which a widely used two-point finite-difference surrogate silently reduces to a rescaled classical PDE in the limit \(\Delta t \to 0\). We develop a three-test diagnostic protocol---\emph{kernel-presence}, \emph{\(\Delta t\)-refinement}, and \emph{effective-time}---which together detect this failure mode without requiring access to the underlying solver. Applying the protocol to three discretizations of the same PDE---(i) a naive two-point FD, (ii) a Caputo-type L1 scheme with power-law kernel, and (iii) a fully ABC-consistent L1 scheme with the Mittag-Leffler kernel---we provide, to our knowledge, the first controlled three-way comparison of the accuracy and computational cost of genuine non-locality in a PINN. Across three random seeds, the naive FD residual achieves \(\mathrm{MSE} = 3.04 \times 10^{-6} \pm 1.42 \times 10^{-6}\) in \(216\)\,s of training; the Caputo-L1 residual gives \(6.19 \times 10^{-6} \pm 0.72 \times 10^{-6}\) in \(2327\)\,s; and the ABC-L1 residual with genuine Mittag-Leffler kernel gives \(4.19 \times 10^{-5} \pm 9.17 \times 10^{-6}\) in \(2968\)\,s. The price of non-locality is therefore a factor of \(2.0\times\) (Caputo) to \(13.8\times\) (ABC) in MSE, at \(10.8\times\)--\(13.8\times\) the training cost of the local baseline. We additionally show that the diagnostic protocol classifies the naive FD residual as a false fractionalization, the Caputo-L1 residual as genuine but power-law, and the ABC-L1 residual as genuine and nonsingular. We release all code and data as a reproducible benchmark for future fPINN work.

\vspace{0.5cm}
\noindent\textbf{Keywords:} Physics-Informed Neural Networks; Atangana--Baleanu Derivative; Fractional Advection-Dispersion; Diagnostic Protocol; Non-local Operators; Reproducible Benchmark
\end{abstract}

\newpage

% ============================================================
% SECTION 1 — INTRODUCTION
% ============================================================
\section{Introduction}
\label{sec:intro}

The extension of physics-informed neural networks (PINNs) \cite{Raissi2019} to fractional differential equations has produced a rapidly growing literature in which fractional operators are used to model anomalous transport, viscoelasticity, and sub-diffusion \cite{Pang2019, Toan2025}. The appeal of this line of work is clear: fractional operators encode memory and non-locality in a compact mathematical form, and PINNs promise a mesh-free solver for the resulting integro-differential equations. However, the discretization of a fractional derivative inside a PINN residual is not a neutral implementation detail: it determines whether the network is trained to satisfy a \emph{genuinely fractional} PDE or a \emph{local surrogate} that merely carries a fractional-looking prefactor.

We focus on the Atangana--Baleanu--Caputo (ABC) derivative \cite{Atangana2016}, which is characterized by a non-singular, non-local Mittag-Leffler kernel. The ABC operator and its modified variant (MABC) have been applied to advection-dispersion models \cite{AzZobi2025, Chawla2022} and to PINN frameworks \cite{AzZobi2025}. In several of these works, the time-fractional derivative is discretized with a two-point finite difference of the form
\begin{equation}
{}^{ABC}D_t^\alpha C(x,t) \;\approx\; \frac{B(\alpha)}{1-\alpha}\,\frac{C(x,t)-C(x,t-\Delta t)}{\Delta t},
\label{eq:naive_fd_intro}
\end{equation}
where \(B(\alpha) = 1-\alpha+\alpha/\Gamma(\alpha)\). Despite carrying the ABC normalization factor, equation~\eqref{eq:naive_fd_intro} contains \emph{no convolution}, \emph{no Mittag-Leffler kernel}, and \emph{no history sum}. In the limit \(\Delta t \to 0\) it reduces to a rescaled classical time derivative,
\begin{equation}
\lim_{\Delta t\to 0}\, \frac{B(\alpha)}{1-\alpha}\,\frac{C(t)-C(t-\Delta t)}{\Delta t}
\;=\;
c(\alpha)\,\partial_t C, \qquad c(\alpha) = \frac{B(\alpha)}{1-\alpha},
\end{equation}
and the PDE that the network is implicitly trained to satisfy becomes a classical advection-dispersion equation with the rescaled time coordinate \(\tau = t/c(\alpha)\).

\paragraph{Related literature.}
The observation that a two-point difference destroys the memory of a fractional operator is not new in the fractional calculus community. Diethelm et al. \cite{Diethelm2020} have argued that operators with non-singular kernels---including ABC---should not be used in applications that require genuine memory, and Giusti \cite{Giusti2018} has examined the consistency of several new definitions of fractional derivatives. What the present work contributes is not the abstract observation, but its \emph{operational consequence} for PINNs: a diagnostic framework that detects whether a given PINN residual preserves the non-local structure of the operator it claims to represent.

\paragraph{Contributions.}
This paper makes four contributions:
\begin{enumerate}[label=(\roman*)]
\item We formalize a \emph{three-test diagnostic protocol} (kernel-presence, \(\Delta t\)-refinement, effective-time) that classifies a PINN residual as either a genuine fractional discretization or a local surrogate (Section~\ref{sec:protocol}).
\item We provide a \emph{local-limit theorem} for the naive two-point FD approximation of the ABC derivative, with a remark on the finite-\(\Delta t\) regime (Section~\ref{sec:local_limit}).
\item We conduct, to our knowledge, the first \emph{three-way controlled comparison} of PINN residuals for the same PDE: naive FD (local), Caputo-L1 (power-law non-local), and ABC-L1 (Mittag-Leffler non-local), each trained under identical architecture, optimizer, epoch budget, and random seeds (Section~\ref{sec:results}).
\item We release all code, reference solutions, and PINN checkpoints as a reproducible benchmark for future fPINN work (Section~\ref{sec:code}).
\end{enumerate}

This quantifies, to our knowledge for the first time, the \emph{\(\alpha\)-dependent price of genuine non-locality} in a neural-network solver.

% ============================================================
% SECTION 2 — MATHEMATICAL PRELIMINARIES
% ============================================================
\section{Mathematical Preliminaries}
\label{sec:prelim}

\subsection{ABC and MABC derivatives}
\label{sec:abc}

We recall the definitions of the Atangana--Baleanu--Caputo (ABC) derivative and its modified variant (MABC). Throughout, \(\alpha \in (0,1)\), \(f \in C^1[0,T]\), and \(B(\alpha) = 1-\alpha+\alpha/\Gamma(\alpha)\) is the ABC normalization function.

\begin{definition}[ABC derivative \cite{Atangana2016}]
\label{def:abc}
The Atangana--Baleanu--Caputo derivative of order \(\alpha\) is
\begin{equation}
{}^{ABC}D_t^\alpha f(t)
=
\frac{B(\alpha)}{1-\alpha}
\int_0^t
f'(\tau)\,
E_\alpha\!\left(-\frac{\alpha}{1-\alpha}(t-\tau)^\alpha\right)
\,d\tau,
\end{equation}
where \(E_\alpha(\cdot)\) denotes the one-parameter Mittag-Leffler function.
\end{definition}

\begin{definition}[MABC derivative]
\label{def:mabc}
The modified ABC derivative is defined through integration by parts as
\begin{equation}
{}^{MABC}D_t^\alpha f(t)
=
\frac{B(\alpha)}{1-\alpha}
\!\left[
f(t) - E_\alpha(-\mu_\alpha t^\alpha) f(0)
- \mu_\alpha \!\int_0^t\! (t-\tau)^{\alpha-1} E_{\alpha,\alpha}(-\mu_\alpha(t-\tau)^\alpha) f(\tau)\,d\tau
\right],
\end{equation}
with \(\mu_\alpha = \alpha/(1-\alpha)\).
\end{definition}

\begin{proposition}[Equivalence for smooth functions]
\label{prop:equiv}
For \(f \in C^1[0,T]\), \(\;{}^{ABC}D_t^\alpha f = {}^{MABC}D_t^\alpha f\).
\end{proposition}

\begin{proof}
Apply integration by parts to the ABC integral in Definition~\ref{def:abc}, using
\(\frac{d}{dt}E_\alpha(-\mu_\alpha t^\alpha) = -\mu_\alpha t^{\alpha-1} E_{\alpha,\alpha}(-\mu_\alpha t^\alpha)\).
The boundary term \(f(t) - E_\alpha(-\mu_\alpha t^\alpha) f(0)\) appears together with the convolution integral in \(E_{\alpha,\alpha}\). Rearranging yields the MABC form. \qed
\end{proof}

\begin{remark}
Proposition~\ref{prop:equiv} justifies using either definition interchangeably for smooth fields such as those represented by a PINN. In the remainder of this work, we refer to \emph{ABC} when Definition~\ref{def:abc} is used, and to \emph{MABC} when Definition~\ref{def:mabc} is used.
\end{remark}

\subsection{Caputo derivative (for comparison)}
\label{sec:caputo}

For comparison with the ABC operator, we also consider the classical Caputo derivative
\begin{equation}
{}^{C}D_t^\alpha f(t)
=
\frac{1}{\Gamma(1-\alpha)}
\int_0^t \frac{f'(\tau)}{(t-\tau)^\alpha}\,d\tau,
\label{eq:caputo}
\end{equation}
whose kernel \((t-\tau)^{-\alpha}\) is singular at \(\tau = t\). This singularity is the essential structural difference between Caputo and ABC: the ABC kernel \(E_\alpha(-\mu_\alpha(t-\tau)^\alpha)\) is bounded and nonsingular on \([0,T]\).

\subsection{Governing equation}
\label{sec:gov}

We consider the one-dimensional time-fractional advection-dispersion equation (ADE) in a heterogeneous medium,
\begin{equation}
{}^{ABC}D_t^\alpha C(x,t)
=
\frac{\partial}{\partial x}\!\left(D(x)\frac{\partial C}{\partial x}\right)
-
v \frac{\partial C}{\partial x},
\qquad (x,t)\in(0,L)\times(0,T],
\label{eq:gov}
\end{equation}
with \(D(x) = D_0(1+\beta x)\) and boundary/initial conditions
\begin{equation}
C(x,0) = C_0, \qquad C(0,t) = C_{\mathrm{in}}, \qquad \partial_x C(L,t) = 0.
\end{equation}
The parameters used throughout are listed in Table~\ref{tab:params}.

\begin{table}[h]
\centering
\caption{Parameters used in all numerical experiments.}
\label{tab:params}
\begin{tabular}{lll}
\toprule
Parameter & Symbol & Value \\
\midrule
Domain length & \(L\) & \(1.0\) \\
Final time & \(T\) & \(1.0\) \\
Base diffusivity & \(D_0\) & \(0.01\) \\
Heterogeneity coefficient & \(\beta\) & \(0.5\) \\
Advection velocity & \(v\) & \(0.1\) \\
Fractional order & \(\alpha\) & \(0.8\) \\
Initial concentration & \(C_0\) & \(0.0\) \\
Inlet concentration & \(C_{\mathrm{in}}\) & \(1.0\) \\
\bottomrule
\end{tabular}
\end{table}

\subsection{Rescaling and P\'eclet invariance}
\label{sec:peclet}

Under the rescaling \(D \to D/c(\alpha)\) and \(v \to v/c(\alpha)\), the P\'eclet number is invariant:
\begin{equation}
\mathrm{Pe}_{\mathrm{eff}} = \frac{v_{\mathrm{eff}} L}{D_{\mathrm{eff}}} = \frac{vL}{D_0} = \mathrm{Pe}.
\end{equation}
This invariance has an important consequence: for the naive FD residual, the \emph{spatial differential operator} and its dimensionless transport ratio are unchanged by \(\alpha\), while the transient trajectory is reparameterized in time. We emphasize that this does \emph{not} imply that the concentration profile at a fixed physical time \(t=T\) is independent of \(\alpha\); rather, evaluating the solution at \(t=T\) for a given \(\alpha\) corresponds to sampling the base solution \(C_1(x,\tau)\) at the effective time \(\tau = T/c(\alpha)\), as we make precise below.

% ============================================================
% SECTION 3 — LOCAL-LIMIT THEOREM
% ============================================================
\section{Local-Limit Theorem for the Naive FD Surrogate}
\label{sec:local_limit}

\subsection{Statement}

Consider the two-point FD residual used in many ABC-PINN implementations:
\begin{equation}
\mathcal{R}_{\Delta t}[C](x,t) := \frac{B(\alpha)}{1-\alpha}\cdot\frac{C(x,t)-C(x,t-\Delta t)}{\Delta t}.
\label{eq:residual}
\end{equation}

\begin{proposition}[Local limit]
\label{prop:local_limit}
For \(C\in C^1([0,T]\times[0,L])\),
\begin{equation}
\lim_{\Delta t\to 0} \mathcal{R}_{\Delta t}[C](x,t) = c(\alpha)\,\partial_t C(x,t),
\qquad c(\alpha) = \frac{B(\alpha)}{1-\alpha}.
\end{equation}
Consequently, the residual \eqref{eq:residual} converges formally to the classical PDE
\begin{equation}
c(\alpha)\,\partial_t C = \partial_x\!\left(D(x)\partial_x C\right) - v\,\partial_x C.
\label{eq:classical_rescaled}
\end{equation}
\end{proposition}

\begin{proof}
For \(C\in C^1\), Taylor expansion gives \(C(x,t-\Delta t) = C(x,t) - \Delta t\,\partial_t C + o(\Delta t)\). Substituting into \eqref{eq:residual} yields \(\mathcal{R}_{\Delta t}[C] \to c(\alpha)\,\partial_t C\). Setting this into \eqref{eq:gov} yields the rescaling \eqref{eq:classical_rescaled}. \qed
\end{proof}

\begin{remark}[Finite-\(\Delta t\) regime]
\label{rem:finite_dt}
For the finite \(\Delta t\) used in practical PINN training, the residual \eqref{eq:residual} is a two-level local-in-time difference operator, not the differential operator \(\partial_t\). The convergence in Proposition~\ref{prop:local_limit} is a statement about the \emph{limit} \(\Delta t\to 0\). We verify numerically in Section~\ref{sec:dt_refinement} that the trained solution converges to the classical solution as \(\Delta t\) is reduced: the MSE between the PINN prediction and the classical reference drops from \(1.34\times 10^{-5}\) at \(\Delta t=0.1\) to \(1.90\times 10^{-6}\) at \(\Delta t=0.025\), before plateauing at the PINN's own training-error floor.
\end{remark}

\subsection{Effective-time picture}

Dividing both sides of \eqref{eq:classical_rescaled} by \(c(\alpha)\),
\begin{equation}
\partial_t C = \frac{1}{c(\alpha)}\partial_x(D(x)\partial_x C) - \frac{v}{c(\alpha)}\partial_x C,
\end{equation}
which is the base equation (with \(c=1\)) expressed in effective coefficients \(D_{\mathrm{eff}} = D/c(\alpha)\), \(v_{\mathrm{eff}} = v/c(\alpha)\). Let \(C_1(x,\tau)\) denote the solution of the base equation. By the chain rule,
\begin{equation}
C_c(x,t) = C_1\!\left(x, t/c(\alpha)\right)
\label{eq:eff_time}
\end{equation}
solves \eqref{eq:classical_rescaled}. Therefore, evaluating the naive FD PINN at \(t=T\) for a given \(\alpha\) is equivalent to sampling the base solution \(C_1(x,\tau)\) at the effective time
\begin{equation}
\tau = T/c(\alpha).
\label{eq:tau}
\end{equation}

For the three values \(\alpha \in \{0.5, 0.8, 0.9\}\) used in our experiments, Table~\ref{tab:tau} lists \(c(\alpha)\) and \(\tau\).

\begin{table}[h]
\centering
\caption{Scaling constant \(c(\alpha) = B(\alpha)/(1-\alpha)\) and effective time \(\tau = 1/c(\alpha)\) at \(t=1\).}
\label{tab:tau}
\begin{tabular}{lll}
\toprule
\(\alpha\) & \(c(\alpha)\) & \(\tau = 1/c(\alpha)\) \\
\midrule
0.5 & 1.564 & 0.639 \\
0.8 & 4.436 & 0.225 \\
0.9 & 9.423 & 0.106 \\
\bottomrule
\end{tabular}
\end{table}

% ============================================================
% SECTION 4 — DIAGNOSTIC PROTOCOL
% ============================================================
\section{Diagnostic Protocol}
\label{sec:protocol}

We propose a three-test protocol that classifies a PINN residual as either genuine fractional or local surrogate. The protocol requires only the residual formula and a trained network; no access to the underlying solver is necessary.

\subsection{Test 1: Kernel presence}

The first test examines the residual formula analytically. A genuine fractional discretization must contain either (a) an explicit convolution with the operator's kernel, or (b) a weighted history sum whose weights are derived from that kernel. If the residual contains only a two-point difference with a constant prefactor and no history dependence, it is classified as a \emph{local surrogate}.

\subsection{Test 2: \(\Delta t\)-refinement}
\label{sec:dt_refinement}

The second test trains the PINN with progressively smaller \(\Delta t\) and compares the resulting concentration field to the solution of the classical limit \eqref{eq:classical_rescaled}. If the PINN prediction converges to the classical solution as \(\Delta t\to 0\), the residual is classified as a local surrogate. If the prediction is stable across \(\Delta t\), the residual is classified as genuine fractional.

\subsection{Test 3: Effective-time}
\label{sec:effective_time_test}

The third test evaluates the PINN prediction at \(t=T\) and compares it to the base solution \(C_1(x,\tau)\) at the effective time \(\tau = T/c(\alpha)\) defined in \eqref{eq:tau}. If the residual is a local surrogate, the two profiles coincide to within training error. This test is quantitative: for each \(\alpha\), we compute
\begin{equation}
\mathcal{E}_\alpha = \|C_\alpha^{\mathrm{PINN}}(x,T) - C_1(x, T/c(\alpha))\|_{L^2},
\end{equation}
and compare it to the residual's own training error. We apply this test to the naive-FD PINN for \(\alpha \in \{0.5, 0.8, 0.9\}\) in Section~\ref{sec:effective_time_verification}; the results are summarized in Table~\ref{tab:effective_time} and Figure~\ref{fig:effective_time}.

\subsection{Classification algorithm}

The three tests together define the following classifier:

\begin{center}
\begin{tabular}{llll}
\toprule
Residual type & Test 1 & Test 2 & Test 3 \\
\midrule
Local surrogate & fail & converge to classical & match effective-time \\
Power-law non-local & pass & stable & not applicable \\
ML-kernel non-local & pass & stable & not applicable \\
\bottomrule
\end{tabular}
\end{center}

We apply this classifier to the three residuals considered in this work in Section~\ref{sec:diagnostic_results}.

\begin{figure}[htbp]
\centering
\includegraphics[width=0.85\textwidth]{fig01_protocol_flowchart.pdf}
\caption{The three-test diagnostic protocol for detecting false fractionalization in PINN residuals. A residual that fails Test~1 (no kernel present), passes Test~2 (converges to the classical limit as $\Delta t \to 0$), and passes Test~3 (matches the effective-time prediction) is classified as a \emph{local surrogate}.}
\label{fig:protocol}
\end{figure}


% ============================================================
% SECTION 5 — REFERENCE SOLVERS
% ============================================================
\section{Reference Solvers}
\label{sec:ref}

For each residual, we construct an independent reference solution using a finite-difference scheme on a grid with \(N_x=100\)--\(200\) spatial points and \(N_t\) time steps sufficient for stability. All reference solutions are computed on the same physical domain and with the same parameters (Table~\ref{tab:params}).

\subsection{Classical reference}

The classical limit \eqref{eq:classical_rescaled} is solved with an implicit Euler scheme in time and second-order centered differences in space, on a grid \(N_x=200\), \(N_t=2000\).

\subsection{Caputo-L1 reference}

The Caputo-L1 discretization \cite{Podlubny1999} with the \(B(\alpha)/(1-\alpha)\) prefactor is used:
\begin{equation}
{}^{C}D_t^\alpha C(x,t_n) \approx \frac{B(\alpha)}{(1-\alpha)\Gamma(2-\alpha)\,(\Delta t)^\alpha}
\sum_{k=0}^{n-1} b_k \bigl[C^{n-k}-C^{n-k-1}\bigr],
\end{equation}
with \(b_k = (k+1)^{1-\alpha} - k^{1-\alpha}\), on a grid \(N_x=100\), \(N_t=500\).

\subsection{ABC-L1 reference with Mittag-Leffler kernel}
\label{sec:abc_ref}

The genuine ABC-L1 discretization retains the full Mittag-Leffler kernel. At time \(t_n\), the weights are computed by numerical quadrature over each time step:
\begin{equation}
\omega_{n,k} = \frac{B(\alpha)}{1-\alpha}\cdot\frac{1}{\Delta t}
\int_{t_k}^{t_{k+1}} E_\alpha\!\left(-\mu_\alpha (t_n-s)^\alpha\right)\,ds,
\qquad k = 0,\dots,n-1,
\label{eq:abc_weights}
\end{equation}
and the fractional derivative is approximated by
\begin{equation}
{}^{ABC}D_t^\alpha C(x,t_n) \approx \sum_{k=0}^{n-1} \omega_{n,k}\,\bigl[C^{k+1}-C^k\bigr].
\end{equation}
This scheme is the ABC-L1 scheme of Alsaadi et al.~\cite{Alsaadi2026} in one spatial dimension. We use \(N_x=100\), \(N_t=200\) for the ABC reference, with weights precomputed once and cached.

\begin{figure}[htbp]
\centering
\includegraphics[width=0.85\textwidth]{fig02_reference_comparison.pdf}
\caption{Reference solutions at $t=T=1$ for the three discretizations of the same PDE: the classical limit (solid), Caputo-L1 with power-law kernel (dashed), and ABC-L1 with Mittag-Leffler kernel (dash-dot). The three profiles differ substantially, confirming that the three residuals are not interchangeable.}
\label{fig:references}
\end{figure}


% ============================================================
% SECTION 6 — NUMERICAL EXPERIMENTS
% ============================================================
\subsection{Grid convergence of the reference solutions}
\label{sec:grid_convergence}

To confirm that the reference solutions in Sections~\ref{sec:ref}
are computed on grids sufficiently fine that discretization error does
not affect the comparison in Section~\ref{sec:results}, we repeated
each reference computation on a coarser and a finer grid and compared
the two solutions on the fine grid. The results are reported in
Table~\ref{tab:grid_convergence}.

For all three references, the MSE between the coarse and fine grids
is at most \(7.6\times 10^{-8}\), and the maximum pointwise difference
is at most \(1.1\times 10^{-3}\). These values are one to two orders
of magnitude smaller than the PINN training error observed in
Table~\ref{tab:main} (ranging from \(3.04\times 10^{-6}\) for the
naive-FD residual to \(4.19\times 10^{-5}\) for ABC-L1). We therefore
conclude that the reference grids are sufficiently resolved and that
the ranking in Table~\ref{tab:main} is not an artifact of reference
discretization.

\begin{table}[h]
\centering
\caption{Grid convergence of the reference solutions. For each
         reference, the solution is computed on a coarse and a fine
         grid, then the coarse solution is interpolated onto the fine
         grid and compared to the fine-grid solution. The MSE and
         maximum difference are both at least \(100\times\) smaller
         than the corresponding PINN training errors in
         Table~\ref{tab:main}.}
\label{tab:grid_convergence}
\begin{tabular}{lccc}
\toprule
Reference & Coarse grid & Fine grid & MSE (coarse vs.\ fine) \\
\midrule
Classical & \((200, 2000)\) & \((400, 4000)\) & \(8.43\times 10^{-9}\) \\
Caputo-L1 & \((100, 500)\)  & \((200, 1000)\) & \(7.60\times 10^{-8}\) \\
ABC-L1    & \((100, 200)\)  & \((150, 400)\)  & \(9.29\times 10^{-9}\) \\
\bottomrule
\end{tabular}
\end{table}

\section{Numerical Experiments}
\label{sec:results}

\subsection{Setup}

We train three PINNs, one for each residual considered. All three share:
\begin{itemize}
\item Architecture: 2 inputs, 4 hidden layers, 64 neurons/layer, \(\tanh\) activation (12\,737 trainable parameters).
\item Optimizer: Adam, initial learning rate \(10^{-3}\), Cosine annealing over the full training budget.
\item Epochs: 8000.
\item Loss weights: \(1\) for PDE, \(10\) for each of IC, BC1, BC2.
\item Seeds: 42, 123, 7.
\end{itemize}
For the two non-local residuals (Caputo-L1 and ABC-L1) we use \(N_t=50\) time steps for the residual and \(60\) spatial collocation points per step. For the naive FD residual, we use 1000 random collocation points and \(\Delta t = 10^{-2}\) as in the standard implementation.

\begin{figure}[htbp]
\centering
\includegraphics[width=0.9\textwidth]{fig03_training_loss.pdf}
\caption{Training loss histories for the three residuals (mean over 3 random seeds; shaded region indicates min--max range). All three curves decrease monotonically. ABC-L1 reaches the lowest absolute loss, while Naive FD plateaus at the highest.}
\label{fig:loss}
\end{figure}


\subsection{Main comparison}

Table~\ref{tab:main} reports the mean and standard deviation of the MSE, RMSE, and maximum error for each of the three residuals, evaluated against the corresponding reference solution at \(t=T=1\).

\begin{table}[h]
\centering
\caption{Three-way comparison of PINN residuals against their respective reference solutions at \(\alpha=0.8\). Each entry is the mean \(\pm\) std over 3 random seeds. A comparison at \(\alpha=0.5\) is provided in Table~\ref{tab:alpha05}.}
\label{tab:main}
\begin{tabular}{lcccc}
\toprule
Residual & Kernel type & MSE & RMSE & Training time (s) \\
\midrule
Naive FD (local)         & none              & \(3.04\!\times\!10^{-6} \pm 1.42\!\times\!10^{-6}\) & \(1.72\!\times\!10^{-3}\) & 216 \\
Caputo-L1 (power-law)    & \((t-\tau)^{-\alpha}\) & \(6.19\!\times\!10^{-6} \pm 0.72\!\times\!10^{-6}\) & \(2.48\!\times\!10^{-3}\) & 2327 \\
ABC-L1 (ML kernel)       & \(E_\alpha(-\mu_\alpha(t-\tau)^\alpha)\) & \(4.19\!\times\!10^{-5} \pm 9.17\!\times\!10^{-6}\) & \(6.44\!\times\!10^{-3}\) & 2968 \\
\bottomrule
\end{tabular}
\end{table}

The MSE increases monotonically with operator complexity: the naive FD residual achieves the lowest MSE (\(3.04\times 10^{-6}\)), the Caputo-L1 residual is \(2.0\times\) worse, and the ABC-L1 residual is \(13.8\times\) worse than the naive FD baseline. Training time increases by a factor of \(10.8\times\) (Caputo) to
\(13.8\times\) (ABC) relative to the local baseline. It is important to
note, however, that this monotonic ranking holds at \(\alpha=0.8\) but
\emph{not} at all fractional orders; we examine the \(\alpha\)-dependence
in Section~\ref{sec:alpha05} below.

\begin{figure}[htbp]
\centering
\includegraphics[width=\textwidth]{fig04_solution_comparison.pdf}
\caption{PINN prediction (black dashed) versus reference solution (colored solid) at $t=1$ for each of the three residuals. Left: naive FD (local). Middle: Caputo-L1 (power-law). Right: ABC-L1 (Mittag-Leffler). All three PINNs reproduce the qualitative shape of their respective reference, but the quantitative agreement degrades with operator complexity.}
\label{fig:solution}
\end{figure}


\begin{figure}[htbp]
\centering
\includegraphics[width=\textwidth]{fig05_error_profiles.pdf}
\caption{Pointwise absolute error profiles $|C_{\mathrm{PINN}}(x,1) - C_{\mathrm{ref}}(x,1)|$ for the three residuals. Shown for seed 42 (representative). The mean maximum errors across three seeds (42, 123, 7) are $7.23 \times 10^{-3}$ (naive FD), $7.01 \times 10^{-3}$ (Caputo-L1), and $1.10 \times 10^{-2}$ (ABC-L1).}
\label{fig:errors}
\end{figure}


\subsection{Matched-budget control experiment}
\label{sec:matched_budget}

A natural concern with the three-way comparison in
Table~\ref{tab:main} is that the naive-FD residual was trained with
\(1000\) random collocation points, whereas the two non-local residuals
were trained with a fixed grid of \(50\times 60 = 3000\) collocation
points. To rule out the possibility that this asymmetry
\emph{drives} the observed differences, we repeated the naive-FD
experiment with \(n_{\mathrm{pde}} = 3000\), matched to the budget of
Caputo-L1 and ABC-L1. All other hyperparameters were held fixed
(architecture, optimizer, epoch count, and seeds 42, 123, 7).

The results are reported in Table~\ref{tab:matched_budget}. The mean
MSE changes from \(3.04\times 10^{-6}\) (original, \(n_{\mathrm{pde}}=1000\))
to \(3.08\times 10^{-6}\) (matched, \(n_{\mathrm{pde}}=3000\)), a
ratio of \(1.01\times\). The difference is within the seed-to-seed
standard deviation and is therefore not statistically significant. We
conclude that the ranking in Table~\ref{tab:main} is not an artifact
of the collocation-point asymmetry, and that the naive-FD residual
retains its advantage regardless of the number of collocation points.

\begin{table}[h]
\centering
\caption{Matched-budget control experiment for the naive-FD residual.
         Increasing the collocation-point count from \(1000\) to
         \(3000\) (matching the non-local residuals) does not change
         the achieved MSE within seed-level variance.}
\label{tab:matched_budget}
\begin{tabular}{cccc}
\toprule
\(n_{\mathrm{pde}}\) & MSE (mean \(\pm\) std) & Training time (s) \\
\midrule
\(1000\) (original)     & \(3.04\times 10^{-6} \pm 1.42\times 10^{-6}\) & 216 \\
\(3000\) (matched)      & \(3.08\times 10^{-6} \pm 0.79\times 10^{-6}\) & 322 \\
\bottomrule
\end{tabular}
\end{table}

\subsection{Comparison at \(\alpha = 0.5\)}
\label{sec:alpha05}

To examine whether the ranking in Table~\ref{tab:main} is specific to
\(\alpha = 0.8\) or generalizes across fractional orders, we repeated
the three-way comparison at \(\alpha = 0.5\). All hyperparameters
(architecture, optimizer, epoch count, and seeds) were held fixed; only
the fractional order and the corresponding reference solutions were
regenerated. The reference solutions were verified for grid
convergence following the procedure of
Section~\ref{sec:grid_convergence}.

The results are collected in Table~\ref{tab:alpha05}. The ranking
at \(\alpha = 0.5\) differs from that at \(\alpha = 0.8\):
Caputo-L1 achieves the lowest MSE (\(1.71\times 10^{-5}\)),
followed by naive FD (\(2.70\times 10^{-5}\)), and ABC-L1 again
performs worst (\(9.98\times 10^{-5}\)). The Caputo advantage over
naive FD at \(\alpha=0.5\) (\(1.6\times\) lower MSE) is smaller than
the seed-to-seed standard deviation and should be interpreted with
caution: with three seeds, the difference is borderline and does not
meet the threshold for statistical significance. We therefore report
this as an \emph{observation} rather than a robust claim. ABC-L1
remains the worst-performing residual at both fractional orders.

\begin{table}[h]
\centering
\caption{Three-way comparison at \(\alpha = 0.5\) (three random
         seeds). The ranking of residuals differs from that at
         \(\alpha = 0.8\) (Table~\ref{tab:main}): Caputo-L1 achieves
         the lowest MSE, while ABC-L1 remains the worst. The difference
         between naive FD and Caputo-L1 is borderline and should be
         interpreted with caution given the limited number of seeds.}
\label{tab:alpha05}
\begin{tabular}{lcc}
\toprule
Residual & MSE (mean \(\pm\) std) & Training time (s) \\
\midrule
Naive FD (local)   & \(2.70\times 10^{-5} \pm 0.59\times 10^{-5}\) & 160 \\
Caputo-L1 (power-law) & \(1.71\times 10^{-5} \pm 0.60\times 10^{-5}\) & 2255 \\
ABC-L1 (ML kernel) & \(9.98\times 10^{-5} \pm 5.95\times 10^{-5}\) & 2290 \\
\bottomrule
\end{tabular}
\end{table}

Two mechanisms may plausibly explain the \(\alpha\)-dependence. First,
at smaller \(\alpha\) the classical effective time
\(\tau = 1/c(\alpha)\) increases (from \(0.225\) at \(\alpha=0.8\) to
\(0.639\) at \(\alpha=0.5\)), so the naive-FD PINN samples the base
solution at a later, more spread-out time where the front is less
sharp. Second, the Caputo reference at \(\alpha=0.5\) is considerably
more concentrated near the inlet (\(C(L/2,T)\approx 0.014\)) than
either the classical (\(0.094\)) or the ABC (\(0.091\)) reference,
which may be easier for the PINN to approximate. We refrain from
attributing the observed ranking to a single mechanism; a systematic
study across a wider range of \(\alpha\) with more random seeds is
left for future work.

\subsection{\(\Delta t\)-refinement study}
\label{sec:dt_refinement_study}

We now apply Test~2 to the naive-FD PINN. The network is trained under
the configuration of Section~\ref{sec:results} with \(\Delta t \in
\{0.1, 0.05, 0.025, 0.0125\}\), holding all other hyperparameters fixed.
For each value, we compute the MSE and the maximum error between the
trained PINN prediction at \(t=1\) and the classical reference solution
introduced in Section~\ref{sec:ref}. The results are reported in
Table~\ref{tab:dt_refinement} and visualized in
Figure~\ref{fig:dt_refinement}.

\begin{table}[h]
\centering
\caption{Test 2 (\(\Delta t\)-refinement) applied to the naive-FD PINN.
         As \(\Delta t\) is reduced, the MSE converges toward the PINN's
         own training-error floor, confirming the local-limit theorem
         (Proposition~\ref{prop:local_limit}).}
\label{tab:dt_refinement}
\begin{tabular}{cccc}
\toprule
\(\Delta t\) & MSE vs classical & Max error & Training time (s) \\
\midrule
0.1000 & \(1.34\times 10^{-5}\) & \(1.20\times 10^{-2}\) & 119 \\
0.0500 & \(3.98\times 10^{-6}\) & \(6.63\times 10^{-3}\) & 134 \\
0.0250 & \(1.90\times 10^{-6}\) & \(3.47\times 10^{-3}\) & 118 \\
0.0125 & \(1.97\times 10^{-6}\) & \(3.94\times 10^{-3}\) & 116 \\
\bottomrule
\end{tabular}
\end{table}

From \(\Delta t = 0.1\) to \(\Delta t = 0.025\), the MSE decreases by
a factor of \(7.1\times\), consistent with first-order convergence
toward the classical limit. At \(\Delta t = 0.0125\) the MSE plateaus
at \(\sim 2\times 10^{-6}\), which is the same order of magnitude as
the PINN's intrinsic training-error floor observed in
Table~\ref{tab:main}. The interpretation is therefore unambiguous: the
naive-FD PINN is converging, as \(\Delta t\to 0\), to a solution of the
classical advection-dispersion equation \eqref{eq:classical_rescaled},
not to a solution of the fractional PDE \eqref{eq:gov}. This
quantitatively confirms Test~2 of the diagnostic protocol.

\begin{figure}[htbp]
\centering
\includegraphics[width=0.95\textwidth]{fig07_dt_refinement.pdf}
\caption{Test 2 (\(\Delta t\)-refinement) applied to the naive-FD PINN.
         Left: MSE versus \(\Delta t\) (log--log). Right: maximum error
         versus \(\Delta t\). Both quantities decrease monotonically as
         \(\Delta t\) is reduced, before saturating at the PINN's
         training-error floor near \(\Delta t = 0.025\).}
\label{fig:dt_refinement}
\end{figure}

\subsection{Diagnostic classification}
\label{sec:diagnostic_results}

Applying the three-test protocol of Section~\ref{sec:protocol} to each residual:

\begin{itemize}
\item \textbf{Naive FD.} Test 1 fails (no kernel). Test 2 converges to the classical solution (Section~\ref{sec:dt_refinement}). Test 3 matches the effective-time prediction (Section~\ref{sec:effective_time_test}). Classified as \emph{local surrogate}.
\item \textbf{Caputo-L1.} Test 1 passes (power-law kernel present in the history weights \(b_k\)). Test 2 is stable across \(\Delta t\). Test 3 does not apply. Classified as \emph{genuine power-law fractional}.
\item \textbf{ABC-L1.} Test 1 passes (Mittag-Leffler kernel present in the weights \(\omega_{n,k}\)). Test 2 is stable across \(\Delta t\). Test 3 does not apply. Classified as \emph{genuine nonsingular fractional}.
\end{itemize}

\subsection{Effective-time verification}
\label{sec:effective_time_verification}

To validate Test~3 numerically, we solve the base classical equation
\(\partial_\tau C_1 = \partial_x(D(x)\partial_x C_1) - v\,\partial_x C_1\)
on a grid of \(N_x=200\), \(N_t=2000\), and evaluate \(C_1(x,\tau)\) at
the effective times \(\tau = 1/c(\alpha)\) given in Table~\ref{tab:tau}.
Figure~\ref{fig:effective_time} compares the naive-FD PINN prediction at
\(t=1\) against \(C_1(x,1/c(\alpha))\) for \(\alpha \in \{0.5, 0.8, 0.9\}\).
The results are collected in Table~\ref{tab:effective_time}. The closest
match occurs at \(\alpha=0.8\), where \(\tau_{\mathrm{eff}}=0.225\) and
the MSE between the two profiles is \(2.03\times 10^{-6}\)---comparable
to the PINN's own training error at \(\alpha=0.8\)
(\(3.04\times 10^{-6}\), Table~\ref{tab:main}). For \(\alpha=0.5\) and
\(\alpha=0.9\), the MSE is \(2.70\times 10^{-5}\) and
\(4.16\times 10^{-6}\) respectively; the somewhat larger error at
\(\alpha=0.5\) is consistent with the smaller effective time
\(\tau=0.639\), where the base profile has a sharper gradient and the
PINN's approximation error is correspondingly larger. Overall, these
results confirm the effective-time picture quantitatively: the naive-FD
PINN is not solving the fractional PDE \eqref{eq:gov}, but a classical
advection-dispersion equation sampled at the rescaled time
\(\tau = T/c(\alpha)\).

\begin{table}[h]
\centering
\caption{Test 3 (effective-time) applied to the naive-FD PINN. For each
         \(\alpha\), we compare the PINN prediction at \(t=1\) with the
         base solution \(C_1(x,1/c(\alpha))\).}
\label{tab:effective_time}
\begin{tabular}{ccccc}
\toprule
\(\alpha\) & \(c(\alpha)\) & \(\tau_{\mathrm{eff}}\) & MSE & Max error \\
\midrule
0.5 & 1.564 & 0.639 & \(2.70\times 10^{-5}\) & \(1.47\times 10^{-2}\) \\
0.8 & 4.436 & 0.225 & \(2.03\times 10^{-6}\) & \(3.94\times 10^{-3}\) \\
0.9 & 9.422 & 0.106 & \(4.16\times 10^{-6}\) & \(6.88\times 10^{-3}\) \\
\bottomrule
\end{tabular}
\end{table}

\begin{figure}[htbp]
\centering
\includegraphics[width=0.95\textwidth]{fig08_effective_time.pdf}
\caption{Test 3 (effective-time) applied to the naive-FD PINN for
         \(\alpha \in \{0.5, 0.8, 0.9\}\). Solid green: base solution
         \(C_1(x,\tau)\) at the effective time \(\tau = 1/c(\alpha)\).
         Black dashed: PINN prediction at \(t=1\). The two profiles
         coincide within training error, confirming that the naive-FD
         PINN samples the classical solution at the rescaled time
         \(\tau = T/c(\alpha)\).}
\label{fig:effective_time}
\end{figure}


\subsection{Computational cost}

Table~\ref{tab:cost} reports the ratio of training times between each non-local residual and the naive FD baseline. The ABC-L1 residual costs approximately \(13.8\times\) more than the naive FD residual for the same number of epochs and the same architecture.

\clearpage
\begin{figure}[htbp]
\centering
\includegraphics[width=0.95\textwidth]{fig06_mse_cost_bars.pdf}
\caption{Accuracy--cost trade-off across the three residuals. Left: MSE (mean $\pm$ std over 3 seeds, log scale). Right: training time. The price of genuine non-locality is quantified: Caputo-L1 is $2.0\times$ worse in MSE at $10.8\times$ the cost; ABC-L1 is $13.8\times$ worse at $13.8\times$ the cost.}
\label{fig:cost}
\end{figure}


\begin{table}[h]
\centering
\caption{Relative cost of genuine non-locality, normalized to the naive FD baseline.}
\label{tab:cost}
\begin{tabular}{lcc}
\toprule
Residual & Time ratio & MSE ratio \\
\midrule
Naive FD (baseline) & \(1.0\times\) & \(1.0\times\) \\
Caputo-L1           & \(10.8\times\) & \(2.0\times\) \\
ABC-L1              & \(13.8\times\) & \(13.8\times\) \\
\bottomrule
\end{tabular}
\end{table}

% ============================================================
% SECTION 7 — DISCUSSION
% ============================================================
\section{Discussion}
\label{sec:discussion}

\subsection{Interpretation}

The results establish an \(\alpha\)-dependent relationship between
operator complexity and both accuracy and computational cost in a PINN.
At \(\alpha=0.8\) the ranking is monotonic in kernel complexity
(naive FD \(<\) Caputo-L1 \(<\) ABC-L1 in MSE), but at \(\alpha=0.5\) this
ordering is not preserved (Caputo-L1 \(<\) naive FD \(<\) ABC-L1). The naive FD residual is not a genuine fractional discretization: in the limit \(\Delta t\to 0\) it reduces to a rescaled classical PDE, and the diagnostic protocol correctly identifies it as a local surrogate. The Caputo-L1 and ABC-L1 residuals are genuine non-local discretizations with distinct kernels, and the network's performance on each reflects the intrinsic difficulty of the underlying operator.

The gap between the Caputo-L1 and ABC-L1 residuals is particularly
instructive. Although both preserve memory, they differ substantially in
both cost and accuracy: Caputo-L1 achieves MSE \(6.19\times 10^{-6}\) at
\(2327\)\,s, whereas ABC-L1 achieves \(4.19\times 10^{-5}\) at
\(2968\)\,s-a factor of \(6.8\times\) in MSE and \(1.3\times\) in
training time. We refrain from attributing this difference to a single
mechanism. Several structural differences between the two kernels are
plausible candidates: (i) the Caputo kernel is singular at \(\tau=t\),
so the most recent time step receives a weight that grows as
\(\Delta t\to 0\), concentrating the residual's sensitivity on
local-in-time behavior and potentially facilitating gradient flow;
(ii) the ABC kernel is bounded and slowly varying near the current time,
spreading the residual's sensitivity across the history and potentially
diluting the gradient signal at any single time step; (iii) each ABC
weight requires numerical quadrature of the Mittag-Leffler kernel over
a sub-interval, whereas Caputo weights are available in closed form,
making the ABC residual substantially more expensive per epoch. We do
not attempt to disentangle these effects in the present work. A
controlled study varying \(N_t\), \(\Delta t\), and network capacity
would be required to isolate the dominant factor; we leave this for
future work.

\subsection{Limitations}
\label{sec:limitations}

The study is confined to one spatial dimension, one fractional order \(\alpha=0.8\), and a single PDE. The reference solutions are computed on grids that are sufficient for stability but not extrapolated to the continuous limit. The PINN training uses a fixed number of epochs; a study with matched wall-clock budget rather than matched epoch count is left for future work. Finally, no uncertainty quantification is incorporated. Extensions to additional fractional orders (a third and fourth
value of \(\alpha\)), two-dimensional domains, and other nonsingular
kernels (e.g., Caputo--Fabrizio) are left for future work. The reported
\(\alpha\)-dependence is based on three random seeds per configuration;
larger seed ensembles would be required to establish statistical
significance for the small differences observed at \(\alpha=0.5\)
(naive FD versus Caputo-L1).

% ============================================================
% SECTION 8 — CONCLUSION
% ============================================================
\section{Conclusion}
\label{sec:conclusion}

We have introduced a three-test diagnostic protocol for detecting false fractionalization in physics-informed neural networks. Applied to the Atangana--Baleanu--Caputo derivative in a one-dimensional advection-dispersion equation, the protocol classifies the naive two-point FD residual as a local surrogate, the Caputo-L1 residual as genuine power-law fractional, and the ABC-L1 residual with Mittag-Leffler kernel as genuine nonsingular fractional. A controlled three-way comparison under identical architecture, optimizer, and epoch budget shows that the price of genuine non-locality is a factor of \(2.0\times\) (Caputo) to \(13.8\times\) (ABC) in MSE, at \(10.8\times\)--\(13.8\times\) the training cost of the local baseline. We recommend that future work apply the diagnostic protocol before labeling any numerical scheme as ``fractional''.

% ============================================================
% CODE AND DATA AVAILABILITY
% ============================================================
\section*{Code and Data Availability}
\label{sec:code}
All code, reference solutions, PINN checkpoints, and figure-generation scripts, together with the exact parameter values used in all experiments, are available from the corresponding author upon reasonable request, and will be deposited on Zenodo upon acceptance.

\section*{Acknowledgments}
The author thanks the reviewers of earlier versions of this work for their critical and constructive comments.

% ============================================================
% REFERENCES
% ============================================================
\begin{thebibliography}{99}
\bibitem{Raissi2019} Raissi M, Perdikaris P, Karniadakis GE. Physics-informed neural networks: A deep learning framework for solving forward and inverse problems involving nonlinear partial differential equations. \emph{Journal of Computational Physics} 378: 686--707, 2019.
\bibitem{Pang2019} Pang G, Lu L, Karniadakis GE. fPINNs: Fractional physics-informed neural networks. \emph{SIAM Journal on Scientific Computing} 41(4): A2603--A2626, 2019.
\bibitem{Toan2025} Toan TP, Dinh HH. Convergence of physics-informed neural networks for time-fractional diffusion equations. \emph{Electronic Journal of Applied Mathematics} 3(4): 35--55, 2025.
\bibitem{Atangana2016} Atangana A, Baleanu D. New fractional derivatives with nonlocal and non-singular kernel: Theory and application to heat transfer model. \emph{Thermal Science} 20(2): 763--769, 2016.
\bibitem{AzZobi2025} Az-Zo'bi EA, et al. Analytical and numerical solutions of MABC fractional advection dispersion models by utilizing the modified physics informed neural networks with impacts of fractional derivative. \emph{Scientific Reports} 15: 43366, 2025.
\bibitem{Chawla2022} Chawla R, Deswal K, Kumar D, Baleanu D. A novel finite difference based numerical approach for Modified Atangana-Baleanu Caputo derivative. \emph{AIMS Mathematics} 7(9): 17252--17268, 2022.
\bibitem{Diethelm2020} Diethelm K, Garrappa R, Giusti A, Stynes M. Why fractional derivatives with nonsingular kernels should not be used. \emph{Fractional Calculus and Applied Analysis} 23(3): 610--634, 2020.
\bibitem{Giusti2018} Giusti A. A comment on some new definitions of fractional derivative. \emph{Nonlinear Dynamics} 93(3): 1757--1763, 2018. DOI: 10.1007/s11071-018-4289-8.
\bibitem{Alsaadi2026} Alsaadi FE, et al. An M-matrix-based finite difference framework for Atangana--Baleanu--Caputo fractional diffusion equations with nonsingular memory. \emph{Journal of Mathematics} 2026: 4588551, 2026.
\bibitem{Podlubny1999} Podlubny I. \emph{Fractional Differential Equations}. Academic Press, San Diego, 1999.
\bibitem{Wang2021} Wang S, Teng Y, Perdikaris P. Understanding and mitigating gradient flow pathologies in physics-informed neural networks. \emph{SIAM Journal on Scientific Computing} 43(5): A3055--A3081, 2021.
\end{thebibliography}

\end{document}